\uuid{6bT4}
\titre{ Décomposition d’un sous-espace de $\mathbb{R}^3$ }

\niveau{}
\module{ Espaces vectoriels }
\chapitre{Base d'un espace vectoriel}
\sousChapitre{}
\theme{}
\auteur{Quentin Liard}
\datecreate{ 2026-02-19}
\organisation{AMSCC}
\difficulte{2}

\contenu{

\texte{
On considère dans $\mathbb{R}^3$ l’ensemble
\[
E=\{(x,y,z)\in\mathbb{R}^3 \mid x+y+z=0\}.
\]
On pose également $
F=\mathrm{Vect}\big((1,-1,0)\big).
$
}

\begin{enumerate}

\item   \question{Montrer que $E$ est un sous-espace vectoriel de $\mathbb{R}^3$.}
\indication{}
\reponse{$E$ est l’ensemble des solutions de l’équation linéaire homogène $x+y+z=0$. Il contient le vecteur nul, est stable par addition et multiplication par un scalaire. Donc $E$ est un sous-espace vectoriel.}

\item   \question{Déterminer $\dim E$ et donner une base de $E$.}
\indication{}
\reponse{On a $z=-x-y$. Donc
\[
(x,y,z)=x(1,0,-1)+y(0,1,-1).
\]
Ainsi $E=\mathrm{Vect}\big((1,0,-1),(0,1,-1)\big)$.
Ces deux vecteurs sont libres, donc $\dim E=2$.}

\item   \question{Montrer que $F$ est inclus dans $E$.}
\indication{}
\reponse{Pour le vecteur $(1,-1,0)$, on a $1-1+0=0$, donc il appartient à $E$. Comme $F$ est l’espace engendré par ce vecteur, on a $F\subset E$.}

\item   \question{Trouver un sous-espace $G$ de dimension 1 tel que $E=F\oplus G$.}
\indication{}
\reponse{On peut prendre un vecteur de $E$ non colinéaire à $(1,-1,0)$, par exemple $(1,0,-1)$.
On pose $G=\mathrm{Vect}((1,0,-1))$.
Les vecteurs $(1,-1,0)$ et $(1,0,-1)$ sont libres et engendrent $E$, donc
\[
E=F\oplus G.
\]
}

\item   \question{L’espace $E$ est-il supplémentaire de $\mathrm{Vect}((1,1,1))$ dans $\mathbb{R}^3$ ?}
\indication{}
\reponse{Oui. On a $\dim E=2$ et $\dim\mathrm{Vect}((1,1,1))=1$.
De plus, si $(x,y,z)\in E\cap\mathrm{Vect}((1,1,1))$, alors $(x,y,z)=\lambda(1,1,1)$ et $x+y+z=3\lambda=0$, donc $\lambda=0$.
Ainsi l’intersection est réduite à $\{0\}$ et
\[
\mathbb{R}^3=E\oplus \mathrm{Vect}((1,1,1)).
\]
}

\end{enumerate}

}